\section{Deterministic Multi-Vector Methodology}
\label{sec:method}

This section formalizes the mathematical and algorithmic foundation of SCDS. The system evaluates each candidate dependency $d = (p, v, e)$---comprising package name $p$, declared version $v$, and ecosystem $e \in \{\text{npm}, \text{PyPI}, \text{Maven}\}$---across three orthogonal threat vectors, culminating in a deterministic risk score $S_{\text{final}} \in [0.0, 10.0]$.

\begin{figure*}[t]
\centering
\resizebox{0.95\textwidth}{!}{% Architecture Diagram for SCDS
\begin{tikzpicture}[
    node distance=1.1cm and 0.6cm,
    box/.style={rectangle, draw=black!75, fill=gray!8, thick, rounded corners=2pt, minimum height=0.7cm, align=center, font=\small},
    secbox/.style={rectangle, draw=black!85, fill=blue!8, thick, rounded corners=2pt, minimum height=0.75cm, align=center, font=\small},
    arrow/.style={-{Stealth[scale=1.0]}, thick, draw=black!80},
    groupbox/.style={draw=black!40, dashed, rounded corners=4pt, inner sep=6pt}
]

% Manifest Ingestion
\node[box, fill=orange!10] (manifest) {\textbf{Dependency Manifest}\\(\texttt{package.json} / \texttt{requirements.txt} / \texttt{pom.xml})};
\node[box, below=0.6cm of manifest] (parser) {\textbf{Multi-Ecosystem Parsing Layer}\\(Semver stripping, version bounds, comment filters)};

\draw[arrow] (manifest) -- (parser);

% Analyzers in parallel
\node[secbox, below left=1.1cm and -0.4cm of parser, text width=2.8cm] (typo) {
    \textbf{Typosquatting Analyzer}\\
    {\scriptsize Levenshtein / Jaro-Winkler}\\
    {\scriptsize Normalization Confusion}\\
    \textbf{Weight: 35\%}
};

\node[secbox, below=1.1cm of parser, text width=2.8cm] (vuln) {
    \textbf{Vulnerability Analyzer}\\
    {\scriptsize Google OSV.dev REST API}\\
    {\scriptsize Compound CVSS Scoring}\\
    \textbf{Weight: 45\%}
};

\node[secbox, below right=1.1cm and -0.4cm of parser, text width=2.8cm] (maint) {
    \textbf{Maintainer Analyzer}\\
    {\scriptsize Registry Age \& Dormancy}\\
    {\scriptsize Version Pinning Heuristics}\\
    \textbf{Weight: 20\%}
};

\draw[arrow] (parser) -| (typo.north);
\draw[arrow] (parser) -- (vuln.north);
\draw[arrow] (parser) -| (maint.north);

% Cache sidecar
\node[box, fill=yellow!15, right=0.5cm of vuln, text width=2.1cm, align=center] (cache) {
    \textbf{SQLite Cache}\\
    {\scriptsize WAL Mode}\\
    {\scriptsize 24h Lazy TTL}
};
\draw[<->, dashed, thick] (vuln) -- (cache);
\draw[<->, dashed, thick] (maint.east) -- ++(0.2,0) |- (cache.south);

% Engine and Override
\node[box, fill=purple!10, below=1.0cm of vuln, text width=7.2cm] (engine) {
    \textbf{Deterministic Risk Engine \& Critical Threat Override}\\
    $S_{raw} = 0.45 S_{vuln} + 0.35 S_{typo} + 0.20 S_{maint}$\\
    {\scriptsize Non-linear override prevents dilution if $\max(S_{vuln}, S_{typo}) \ge 7.5$}
};

\draw[arrow] (typo.south) |- (engine.west);
\draw[arrow] (vuln.south) -- (engine.north);
\draw[arrow] (maint.south) |- (engine.east);

% Outputs
\node[box, fill=green!12, below left=0.8cm and 0.2cm of engine, text width=3.8cm] (cli) {
    \textbf{Developer CLI \& CI/CD Gate}\\
    {\scriptsize Rich terminal tables, exit codes (0/1/2)}\\
    {\scriptsize Automated build-breaker}
};

\node[box, fill=cyan!10, below right=0.8cm and 0.2cm of engine, text width=3.8cm] (dash) {
    \textbf{Orchestrator \& Dashboard}\\
    {\scriptsize Recursive workspace discovery}\\
    {\scriptsize FastAPI interactive risk visualizer}
};

\draw[arrow] (engine.south) -| (cli.north);
\draw[arrow] (engine.south) -| (dash.north);

\end{tikzpicture}
}
\caption{Architectural overview of the SCDS pipeline, illustrating the parallel execution of the three analyzer vectors, rate-limit-resistant caching, and the Critical Threat Override Shield.}
\label{fig:arch}
\end{figure*}

\subsection{Vector A: Typosquatting and Brand-Jacking Analysis}
\label{sec:method_typo}
Attackers exploit typographical errors and visual ambiguity to entice developers into installing malicious packages \cite{davis2025typosmart}. The typosquatting analyzer compares candidate name $p$ against a curated ground-truth catalog $\mathcal{C}_e$ of high-reputation packages within ecosystem $e$.

\subsubsection{Normalization Confusion}
In package registries, punctuation variation is a primary confusion mechanism. SCDS first normalizes separators by mapping all hyphens to underscores:
\begin{equation}
\mathcal{N}(s) = \text{replace}(s, \text{'-'}, \text{'\_'})
\end{equation}
If $\mathcal{N}(p) = \mathcal{N}(q)$ for some $q \in \mathcal{C}_e$ while $p \neq q$ (e.g., \texttt{react\_dom} targeting \texttt{react-dom}), the candidate is immediately classified as a high-confidence dependency confusion attack, assigning $S_{\text{typo}} = 9.5$.

\subsubsection{String Distance Formulation}
When no direct separator variant exists, SCDS computes two complementary distance metrics:
\begin{enumerate}
    \item \textbf{Levenshtein Edit Distance \cite{levenshtein1966binary}:} Let $\text{Lev}(s_1, s_2)$ denote the minimum number of single-character insertions, deletions, or substitutions required to transform $s_1$ into $s_2$. SCDS implements this dynamically using two alternating row buffers, ensuring an $O(\min(|s_1|, |s_2|))$ spatial memory footprint.
    \item \textbf{Jaro-Winkler Similarity \cite{winkler1990string}:} Jaro similarity $d_j \in [0.0, 1.0]$ is defined over character matches $m$ within match window $\Delta = \lfloor \max(|s_1|, |s_2|) / 2 \rfloor - 1$ and half-transpositions $t$:
    \begin{equation}
    d_j = \frac{1}{3}\left(\frac{m}{|s_1|} + \frac{m}{|s_2|} + \frac{m - t}{m}\right)
    \end{equation}
    The Jaro-Winkler metric applies a scaling bonus for common prefix length $L \le 4$ with prefix factor $p = 0.1$:
    \begin{equation}
    d_{jw} = \begin{cases} d_j & \text{if } d_j < 0.70 \\ d_j + (L \cdot p \cdot (1 - d_j)) & \text{if } d_j \ge 0.70 \end{cases}
    \end{equation}
\end{enumerate}

\subsubsection{Typosquatting Risk Assignment}
Let $q^* = \arg\max_{q \in \mathcal{C}_e} d_{jw}(p, q)$ be the closest legitimate neighbor, with minimum edit distance $\text{Lev}^* = \text{Lev}(p, q^*)$ and peak similarity $d_{jw}^* = d_{jw}(p, q^*)$. SCDS determines $S_{\text{typo}}$ deterministically:
\begin{equation}
S_{\text{typo}} = \begin{cases}
0.0 & \text{if } p \in \mathcal{C}_e \\
9.5 & \text{if } \text{Lev}^* = 1 \land d_{jw}^* \ge 0.88 \\
7.5 & \text{if } \text{Lev}^* = 2 \land d_{jw}^* \ge 0.85 \\
6.0 & \text{if } d_{jw}^* \ge 0.90 \\
0.0 & \text{otherwise}
\end{cases}
\end{equation}

\subsection{Vector B: Vulnerability Mapping and Compound CVSS}
\label{sec:method_vuln}
Candidate package $(p, v)$ is queried against the Google OSV.dev database \cite{osv2024database}. OSV aggregates CVEs and GHSAs across all supported ecosystems.

\subsubsection{CVSS Parsing}
For each advisory $k \in \mathcal{A}_{(p, v)}$, the analyzer extracts numeric Common Vulnerability Scoring System (CVSS) scores from the JSON payload:
\begin{itemize}
    \item Direct numeric CVSS 3.x vector ratings: $\text{CVSS}_k \in [0.0, 10.0]$.
    \item When only categorical severity strings are present, SCDS maps them monotonically: \texttt{CRITICAL} $\to 9.5$, \texttt{HIGH} $\to 8.0$, \texttt{MEDIUM} $\to 5.5$, and \texttt{LOW} $\to 3.0$.
\end{itemize}

\subsubsection{Compound Aggregation Formula}
Reporting only the single peak CVSS score ignores the compound threat surface of multiple concurrent vulnerabilities, while linear summation rapidly saturates at $10.0$ for benign packages with minor bugs. SCDS applies an asymptotic compound formulation:
\begin{equation}
\begin{split}
S_{\text{vuln}} = \min \Bigl( 10.0, \, & \max_{k \in \mathcal{A}} (\text{CVSS}_k) \\
& + \min\left(1.0, 0.2 \times (|\mathcal{A}| - 1)\right) \Bigr)
\end{split}
\end{equation}
where $|\mathcal{A}|$ is the number of distinct reported advisories. This bounds the compound increment to at most $+1.0$, preserving peak severity dominance while capturing multi-exploit surfaces.

\subsection{Vector C: Maintainer Trust and Metadata Hygiene}
\label{sec:method_maint}
Account compromise and supply chain hygiene deterioration are captured via four empirical signals evaluated against package registry metadata (PyPI JSON and npm registry endpoints):
\begin{equation}
S_{\text{maint}} = \min\left(10.0, f_{\text{pin}} + f_{\text{age}} + f_{\text{dorm}} + f_{\text{repo}}\right)
\end{equation}
where the individual penalty components are defined as:
\begin{itemize}
    \item \textbf{Unpinned Version Penalty ($f_{\text{pin}}$):}
    \begin{equation}
    f_{\text{pin}} = \begin{cases} 2.0 & \text{if } v \text{ is unpinned or wildcard} \\ 0.0 & \text{otherwise} \end{cases}
    \end{equation}
    \item \textbf{Brand-New Package Vector ($f_{\text{age}}$):}
    \begin{equation}
    f_{\text{age}} = \begin{cases} 5.0 & \text{if age } < 14 \text{ days} \\ 2.5 & \text{if } 14 \le \text{age} < 45 \text{ days} \\ 0.0 & \text{otherwise} \end{cases}
    \end{equation}
    \item \textbf{Maintainer Abandonment ($f_{\text{dorm}}$):}
    \begin{equation}
    f_{\text{dorm}} = \begin{cases} 3.5 & \text{if last release } > 1100 \text{ days} \\ 1.5 & \text{if } 730 < \text{last release } \le 1100 \text{ days} \\ 0.0 & \text{otherwise} \end{cases}
    \end{equation}
    \item \textbf{Missing Repository ($f_{\text{repo}}$):} If no repository URL is reported, $f_{\text{repo}} = 2.0$.
\end{itemize}

\subsection{Raw Weighted Integration}
The baseline multi-vector score aggregates the three components using empirically tuned weights that sum to unity:
\begin{equation}
S_{\text{raw}} = (0.45 \times S_{\text{vuln}}) + (0.35 \times S_{\text{typo}}) + (0.20 \times S_{\text{maint}})
\end{equation}
The vulnerability vector is assigned the largest weight ($45\%$) due to direct exploitability, followed by typosquatting ($35\%$) and maintainer hygiene ($20\%$).

\subsection{The Critical Threat Override Shield}
\label{sec:method_override}
A fatal flaw of linear weighted scoring in cybersecurity is \textit{benign dilution}. Consider a package with a confirmed critical remote code execution flaw ($\text{CVSS} = 9.8 \implies S_{\text{vuln}} = 9.8$), zero typosquatting signatures ($S_{\text{typo}} = 0.0$), and pristine maintainer hygiene ($S_{\text{maint}} = 0.0$). Standard linear calculation yields:
\begin{equation}
S_{\text{raw}} = (0.45 \times 9.8) + (0.35 \times 0.0) + (0.20 \times 0.0) = 4.41
\end{equation}
Under standard severity tiers, a score of $4.41$ is categorized as ``Low Risk'' (or ``Moderate''), allowing an unpatched RCE exploit to pass CI/CD build gates completely undetected.

To eliminate this vulnerability, SCDS formalizes the \textbf{Critical Threat Override Shield}:
\begin{equation}
M = \max(S_{\text{vuln}}, S_{\text{typo}})
\end{equation}
\begin{equation}
S_{\text{final}} = \begin{cases}
\max(S_{\text{raw}}, M \times 0.95) & \text{if } M \ge 9.0 \\
\max(S_{\text{raw}}, M \times 0.90) & \text{if } 7.5 \le M < 9.0 \\
S_{\text{raw}} & \text{otherwise}
\end{cases}
\end{equation}
With the override shield active, the candidate with $S_{\text{vuln}} = 9.8$ evaluates to:
\begin{equation}
S_{\text{final}} = \max(4.41, 9.8 \times 0.95) = 9.31 \implies \textbf{CRITICAL}
\end{equation}
This mathematically guarantees that acute threat vectors are never diluted by benign metadata.

\subsection{Risk Tiers and Policy Gate Decisions}
Final scores are rounded to one decimal place and binned into five standard severity tiers:
\begin{itemize}
    \item \textbf{SAFE:} $[0.0, 2.9]$; \textbf{LOW:} $[3.0, 5.4]$; \textbf{MEDIUM:} $[5.5, 7.4]$; \textbf{HIGH:} $[7.5, 8.9]$; \textbf{CRITICAL:} $[9.0, 10.0]$
\end{itemize}
Downstream CI/CD gates evaluate user policies (e.g., \texttt{--fail-on HIGH}), returning standard UNIX process exit codes:
\begin{equation}
\text{ExitCode} = \begin{cases}
0 & \text{if } \forall d \in \mathcal{D}, \, \text{Rank}(d) < \text{Rank}_{\text{policy}} \\
1 & \text{if } \exists d \in \mathcal{D}, \, \text{Rank}(d) \ge \text{Rank}_{\text{policy}} \\
2 & \text{if operational / syntax failure}
\end{cases}
\end{equation}
